\documentclass[11pt,a4paper]{article}
\usepackage[utf8]{inputenc}
\usepackage[spanish]{babel}
\usepackage{amsmath, amssymb}
\usepackage{geometry}
\geometry{top=2.5cm, bottom=2.5cm, left=2.5cm, right=2.5cm}

\begin{document}

\section*{Piense, dialogue y comparta}

\begin{itemize}
    \item[\textbf{1.}] Usted es miembro de un grupo de testigos expertos que presta sus servicios a la comunidad legal. Un abogado defensor le ha pedido a su grupo argumentar en el juicio de un conductor que no excedió el límite de velocidad. Se le proporcionan los siguientes datos: la masa del auto es de $1.50 \times 10^3 \text{ kg}$. La masa del conductor es de $95.0 \text{ kg}$. El coeficiente de fricción cinética entre los neumáticos del auto y la carretera es $0.580$. El coeficiente de fricción estática entre los neumáticos del auto y la carretera es $0.820$. El límite de velocidad registrado en la carretera es de $25 \text{ mi/h}$. En el momento del incidente, la carretera estaba seca y el clima era soleado. También se le proporciona la siguiente descripción del incidente: el conductor manejaba por una colina que hace un ángulo de $17.5^\circ$ con la horizontal. Vio a un perro que corrió hacia la calle, frenó y dejó una marca de $17.0 \text{ m}$ de la longitud que patinó. El auto se detuvo al final de la marca de la patinada. El conductor no le pegó al perro, pero el sonido de las llantas chirriantes atrajo la atención de un policía cercano, que multó al conductor por exceso de velocidad. (a) ¿Debería su grupo estar de acuerdo en ofrecer testimonio en la defensa en este caso? (b) ¿Por qué?
    \vspace*{9cm}

    \item[\textbf{2.}] Considere la actividad de atrapar huevos que se mencionó en el argumento introductorio del capítulo. Discuta en su grupo y haga cálculos de lo siguiente. Identifique una distancia de separación entre el lanzador y el receptor del huevo, y determine una rapidez típica con la que se debe tirar el huevo para que cubra la distancia sin pegar en el suelo o pasar sobre la cabeza del receptor. Estime la masa y el diámetro del huevo (el diámetro más corto, perpendicular a la dimensión más larga). A partir de estos datos, calcule la fuerza sobre el huevo ejercida por su mano si mantiene la mano rígida y no se mueve cuando el huevo la golpea. Ahora, simule mover las manos hacia atrás mientras atrapa el huevo. Que un miembro del grupo estime la distancia que se mueven sus manos en este proceso. A partir de estos datos, estime la fuerza que ejerció la mano sobre el huevo en este proceso de captura. Compare las fuerzas de su mano sobre el huevo entre estos dos métodos para atraparlo.
    \vspace*{8cm}

    \item[\textbf{3.}] \textbf{ACTIVIDAD} Es posible seguir un procedimiento simple para medir el coeficiente de fricción usando la técnica discutida en el ejemplo 5.11. Ponga su libro sobre una mesa y coloque una moneda en la cubierta lejos del lomo. Abra lentamente la cubierta del libro para que forme un plano inclinado hacia abajo por el que la moneda se deslizará. Observe con cuidado y deje de abrir la cubierta en el instante que la moneda comienza a deslizarse. Con un transportador mida este ángulo crítico que la cubierta del libro hace con la horizontal. A partir de este ángulo, determine el coeficiente de fricción estática entre la moneda y la cubierta. (b) Coloque cinta adhesiva entre dos monedas y repita el procedimiento anterior para ambas. ¿Cómo se compara el coeficiente de fricción estática para las dos monedas con el de la única moneda?
    \vspace*{8cm}
\end{itemize}

\newpage

\section*{Problemas}

\subsection*{SECCIÓN 5.1 Concepto de fuerza}

\begin{itemize}
    \item[\textbf{1.}] Cierto ortodoncista utiliza una abrazadera de alambre para alinear la dentadura de un paciente, como se muestra en la figura P5.1. La tensión en el alambre se ajusta para tener $18.0 \text{ N}$. Encuentre la magnitud de la fuerza neta ejercida por el alambre sobre la dentadura.
    \vspace*{6cm}

    \item[\textbf{2.}] Una o más fuerzas externas, suficientemente grandes para ser fácilmente medibles, se ejercen sobre cada objeto encerrado en un recuadro con líneas discontinuas, como se muestra en la figura 5.1. Identifique la reacción a cada una de dichas fuerzas.
    \vspace*{6cm}
\end{itemize}

\subsection*{SECCIÓN 5.4 Segunda ley de Newton}

\begin{itemize}
    \item[\textbf{3.}] Un objeto de $3.00 \text{ kg}$ experimenta una aceleración dada por $\vec{a} = (2.00\hat{i} + 5.00\hat{j})$. Encuentre (a) la fuerza resultante que actúa sobre el objeto y (b) la magnitud de la fuerza resultante.
    \vspace*{6cm}

    \item[\textbf{4.}] La rapidez promedio de una molécula de nitrógeno en el aire es aproximadamente $6.70 \times 10^2 \text{ m/s}$ y su masa es $4.68 \times 10^{-26} \text{ kg}$. (a) Si una molécula de nitrógeno tarda $3.00 \times 10^{-13} \text{ s}$ en golpear una pared y rebotar con la misma rapidez pero moviéndose en la dirección opuesta, ¿cuál es la aceleración promedio de la molécula durante este intervalo de tiempo? (b) ¿Qué fuerza promedio ejerce la molécula sobre la pared?
    \vspace*{6cm}
\end{itemize}

\newpage

\begin{itemize}
    \item[\textbf{5.}] Dos fuerzas, $\vec{F}_1 = (-6.00\hat{i} - 4.00\hat{j})\text{ N}$ y $\vec{F}_2 = (-3.00\hat{i} + 7.00\hat{j})\text{ N}$, actúan sobre una partícula con masa de $2.00 \text{ kg}$ que inicialmente está en reposo en las coordenadas $(-2.00 \text{ m}, 4.00 \text{ m})$. (a) ¿Cuáles son las componentes de la velocidad de la partícula en $t = 10.0 \text{ s}$? (b) ¿En qué dirección se mueve la partícula en el instante $t = 10.0 \text{ s}$? (c) ¿Qué desplazamiento experimenta la partícula durante los primeros $10.0 \text{ s}$? (d) ¿Cuáles son las coordenadas de la partícula en $t = 10.0 \text{ s}$?
    \vspace*{8cm}

    \item[\textbf{6.}] La fuerza ejercida por el viento sobre las sails de un bote es de $390 \text{ N}$ hacia el norte. El agua ejerce una fuerza de $180 \text{ N}$ hacia el este. Si el bote (incluyendo su tripulación) tiene una masa de $270 \text{ kg}$, ¿cuáles son la magnitud y dirección de su aceleración?
    \vspace*{6cm}

    \item[\textbf{7.}] \textbf{Problema de repaso.} Tres fuerzas que actúan sobre un objeto están dadas por $\vec{F}_1 = (-2.00\hat{i} + 2.00\hat{j})\text{ N}$, $\vec{F}_2 = (5.00\hat{i} - 3.00\hat{j})\text{ N}$ y $\vec{F}_3 = (-45.0\hat{i})\text{ N}$. El objeto experimenta una aceleración de $3.75 \text{ m/s}^2$ de magnitud. (a) ¿Cuál es la dirección de la aceleración? (b) ¿Cuál es la masa del objeto? (c) Si el objeto inicialmente está en reposo, ¿cuál es su rapidez después de $10.0 \text{ s}$? (d) ¿Cuáles son las componentes de velocidad del objeto después de $10.0 \text{ s}$?
    \vspace*{8cm}

    \item[\textbf{8.}] Si una sola fuerza constante actúa sobre un objeto que se mueve en una línea recta, la velocidad del objeto es una función lineal del tiempo. La ecuación $v = v_i + at$ da su velocidad $v$ en función del tiempo, donde $a$ es su aceleración constante. ¿Qué pasa si ahora la velocidad es una función lineal de la posición? Suponga que un objeto particular se mueve a través de un medio resistivo, su rapidez disminuye como se describe en la ecuación $v = v_i - kx$, donde $k$ es un coeficiente constante y $x$ es la posición del objeto. Encuentre la ley que describe la fuerza total que actúa sobre este objeto.
    \vspace*{6cm}
\end{itemize}

\newpage

\subsection*{SECCIÓN 5.5 Fuerza gravitacional y peso}

\begin{itemize}
    \item[\textbf{9.}] \textbf{Problema de repaso.} La fuerza gravitacional ejercida sobre una pelota de béisbol es de $2.21 \text{ N}$ hacia abajo. Un pitcher lanza la bola horizontalmente con una velocidad de $18.0 \text{ m/s}$ al acelerarla uniformemente en línea recta horizontal durante un intervalo de tiempo de $170 \text{ ms}$. La pelota parte del reposo. (a) ¿Qué distancia recorre antes de ser liberada? (b) ¿Qué magnitud y dirección tiene la fuerza que el pitcher ejerce sobre la pelota?
    \vspace*{6cm}

    \item[\textbf{10.}] \textbf{Problema de repaso.} La fuerza gravitacional ejercida sobre una pelota de béisbol es $-F_g\hat{j}$. Un pitcher lanza la pelota con velocidad $v\hat{i}$ al acelerarla uniformemente en línea horizontal durante un intervalo de tiempo de $\Delta t = t - 0 = t$. (a) Partiendo del reposo, ¿qué distancia recorre la bola antes de ser liberada? (b) ¿Qué fuerza ejerce el pitcher sobre la pelota?
    \vspace*{6cm}

    \item[\textbf{11.}] \textbf{Problema de repaso.} Un electrón de $9.11 \times 10^{-31} \text{ kg}$ de masa tiene una rapidez inicial de $3.00 \times 10^5 \text{ m/s}$. Viaja en línea recta y su rapidez aumenta a $7.00 \times 10^5 \text{ m/s}$ en una distancia de $5.00 \text{ cm}$. Si supone que su aceleración es constante, (a) determine la fuerza que se ejerce sobre el electrón y (b) compare esta fuerza con el peso del electrón, que se ignoró.
    \vspace*{6cm}

    \item[\textbf{12.}] Si un hombre pesa $900 \text{ N}$ en la Tierra, ¿cuál sería su peso en Júpiter, en donde la aceleración en caída libre es de $25.9 \text{ m/s}^2$?
    \vspace*{5cm}

    \item[\textbf{13.}] Usted está de pie en el asiento de una silla y luego salta. (a) Durante el intervalo de tiempo en el que está en vuelo hacia el suelo, la Tierra se tambalea hacia usted con una aceleración, ¿de qué orden de magnitud? En su solución, explique su lógica. Represente a la Tierra como un objeto perfectamente sólido. (b) La Tierra se mueve hacia arriba a través de una distancia, ¿de qué orden de magnitud?
    \vspace*{6cm}
\end{itemize}

\subsection*{SECCIÓN 5.6 Tercera ley de Newton}

\begin{itemize}
    \item[\textbf{14.}] Un ladrillo de masa $M$ está sobre una almohadilla de hule de masa $m$. Juntos se deslizan hacia la derecha con velocidad constante sobre un estacionamiento cubierto de hielo. (a) Dibuje un diagrama de cuerpo libre del ladrillo e identifique cada fuerza que actúa sobre él. (b) Dibuje un diagrama de cuerpo libre de la almohadilla e identifique cada fuerza que actúa sobre ella. (c) Identifique todos los pares de fuerzas acción-reacción en el sistema ladrillo-almohadilla-planeta.
    \vspace*{7cm}
\end{itemize}

\newpage

\subsection*{SECCIÓN 5.7 Modelos de análisis en los que se utiliza la segunda ley de Newton}

\begin{itemize}
    \item[\textbf{15.}] \textbf{Problema de repaso.} La figura P5.15 muestra un trabajador que empuja un bote, un modo de transporte muy eficiente, a través de un lago tranquilo. Él empuja paralelo a la longitud de la vara ligera y ejerce sobre el fondo del lago una fuerza de $240 \text{ N}$. Suponga que la vara se encuentra en el plano vertical que contiene la quilla del bote. En algún momento, la vara forma un ángulo de $35.0^\circ$ con la vertical y el agua ejerce una fuerza de arrastre horizontal de $47.5 \text{ N}$ sobre el bote, opuesta a su velocidad hacia adelante de $0.857 \text{ m/s}$ de magnitud. La masa del bote, que incluye su carga y al trabajador, es de $370 \text{ kg}$. (a) El agua ejerce una fuerza de flotación vertical hacia arriba sobre el bote. Encuentre la magnitud de esta fuerza. (b) Modele las fuerzas como constantes en un intervalo corto de tiempo para encontrar la velocidad del bote $0.450 \text{ s}$ después del momento descrito.
    \vspace*{6cm}

    \item[\textbf{16.}] Un tornillo de hierro de $65.0 \text{ g}$ de masa cuelga de una cuerda de $35.7 \text{ cm}$ de largo. El extremo superior de la cuerda está fijo. Sin tocarlo, un imán atrae el tornillo de modo que permanece fijo, desplazado horizontalmente $28.0 \text{ cm}$ a la derecha desde la línea vertical previa de la cuerda. El imán está ubicado a la derecha del tornillo y al mismo nivel vertical que el tornillo en la configuración final. (a) Dibuje un diagrama de cuerpo libre del tornillo. (b) Encuentre la tensión en la cuerda. (c) Encuentre la fuerza magnética sobre el tornillo.
    \vspace*{7cm}

    \item[\textbf{17.}] Un bloque se desliza hacia abajo por un plano sin fricción con una inclinación de $\theta = 15.0^\circ$. El bloque parte del reposo en lo alto y la longitud del plano es $2.00 \text{ m}$. (a) Dibuje un diagrama de cuerpo libre del bloque. Encuentre (b) la aceleración del bloque y (c) su rapidez cuando llega al fondo del plano inclinado.
    \vspace*{7cm}

    \item[\textbf{18.}] Un saco de cemento de peso $F_g$ cuelga en equilibrio de tres alambres, como se muestra en la figura P5.18. Dos de los alambres forman ángulos $\theta_1$ y $\theta_2$ con la horizontal. Si supone que el sistema está en equilibrio, demuestre que la tensión en el alambre izquierdo es
    $$ T_1 = \frac{F_g \cos \theta_2}{\operatorname{sen}(\theta_1 + \theta_2)} $$
    \vspace*{6cm}
\end{itemize}

\newpage

\begin{itemize}
    \item[\textbf{19.}] La distancia entre dos postes de teléfono es de $50.0 \text{ m}$. Cuando un ave de $1.00 \text{ kg}$ se posa sobre el alambre telefónico a la mitad entre los postes, el alambre se comba $0.200 \text{ m}$. (a) Dibuje un diagrama de cuerpo libre del ave. (b) ¿Cuánta tensión produce el ave en el alambre? Ignore el peso del alambre.
    \vspace*{6cm}

    \item[\textbf{20.}] Se observa que un objeto de masa $m = 1.00 \text{ kg}$ tiene una aceleración $\vec{a}$ de $10.0 \text{ m/s}^2$ en una dirección a $60.0^\circ$ al noreste, véase la figura P5.20. La fuerza $\vec{F}_2$ que se ejerce sobre el objeto tiene una magnitud de $5.00 \text{ N}$ y se dirige al norte. Determine la magnitud y dirección de la fuerza horizontal $\vec{F}_1$ que actúa sobre el objeto.
    \vspace*{6cm}

    \item[\textbf{21.}] Un acelerómetro simple se construye dentro de un auto al colgar un objeto de masa $m$ de una cuerda de longitud $L$ fija al techo del automóvil. Conforme el auto acelera el sistema cuerda-objeto hace un ángulo constante $\theta$ con la vertical. (a) Si supone que la masa de la cuerda es despreciable comparada con $m$, obtenga una expresión para la aceleración del auto en términos de $\theta$ y demuestre que es independiente de la masa $m$ y la longitud $L$. (b) Determine la aceleración del automóvil cuando $\theta = 23.0^\circ$.
    \vspace*{6cm}

    \item[\textbf{22.}] Un objeto de masa $m_1 = 5.00 \text{ kg}$ colocado sobre una mesa horizontal sin fricción se conecta a una cuerda que pasa sobre una polea y después se une a un objeto colgante de masa $m_2 = 9.00 \text{ kg}$, como se muestra en la figura P5.22. (a) Dibuje diagramas de cuerpo libre de ambos objetos. Encuentre (b) la magnitud de la aceleración de los objetos y (c) la tensión en la cuerda.
    \vspace*{7cm}
\end{itemize}

\newpage

\begin{itemize}
    \item[\textbf{23.}] En el sistema que se muestra en la figura P5.23, una fuerza horizontal $\vec{F}_x$ actúa sobre un objeto de masa $m_2 = 8.00 \text{ kg}$. La superficie horizontal no tiene fricción. Considere la aceleración del objeto deslizante como una función de $F_x$. (a) ¿Para qué valores de $F_x$ el objeto de masa $m_1 = 2.00 \text{ kg}$ acelera hacia arriba? (b) ¿Para qué valores de $F_x$ la tensión en la cuerda es cero? (c) Grafique la aceleración del objeto $m_2$ frente a $F_x$. Incluya valores de $F_x$ desde $-100 \text{ N}$ hasta $+100 \text{ N}$.
    \vspace*{7cm}

    \item[\textbf{24.}] Un auto queda atorado en el lodo. Una grúa jala al automóvil, como se muestra en la figura P5.24. El cable de la grúa tiene una tensión de $2500 \text{ N}$ y jala hacia abajo y hacia la izquierda sobre el perno en su parte superior. El perno ligero se mantiene en equilibrio debido a las fuerzas ejercidas por las barras A y B. Cada barra es un \textit{puntal}; es decir, las barras tienen un peso pequeño comparado con las fuerzas que ellas ejercen en los extremos de los pasadores de bisagras. Cada puntal ejerce una fuerza paralela a su longitud. Determine la fuerza de tensión o de compresión en cada puntal. Proceda como sigue. Haga una conjetura de cómo actúan (empujando o jalando) las fuerzas en el perno superior. Dibuje un diagrama de cuerpo libre del perno. Utilice la condición de equilibrio del perno para trasladar el diagrama de cuerpo libre a ecuaciones. A partir de las ecuaciones calcule las fuerzas ejercidas por los puntales A y B. Si usted obtiene una respuesta positiva, entonces son correctas las direcciones inicialmente propuestas para las fuerzas. Una respuesta negativa significa que la dirección debe invertirse, pero es correcto el valor absoluto de la magnitud de la fuerza. Si un puntal jala sobre el perno, entonces es tensión. Si empuja, entonces el puntal está bajo compresión. Identifique si cada puntal está en tensión o en compresión.
    \vspace*{8cm}

    \item[\textbf{25.}] Un objeto de masa $m_1$ cuelga de una cuerda que pasa por una polea ligera, fija $P_1$, como se muestra en la figura P5.25. La cuerda se conecta a una segunda polea muy ligera $P_2$. Una segunda cuerda pasa por esta polea con un extremo fijo a una pared y el otro a un objeto de masa $m_2$ sobre una mesa horizontal sin fricción. (a) Si $a_1$ y $a_2$ son las aceleraciones de $m_1$ y $m_2$, respectivamente, ¿cuál es la relación entre dichas aceleraciones? Encuentre expresiones para (b) las tensiones en las cuerdas y (c) las aceleraciones $a_1$ y $a_2$ en términos de las masas $m_1$ y $m_2$, y de $g$.
    \vspace*{8cm}
\end{itemize}

\newpage

\subsection*{SECCIÓN 5.8 Fuerzas de fricción}

\begin{itemize}
    \item[\textbf{26.}] ¿Por qué no es posible la siguiente situación? Su libro de física de $3.80 \text{ kg}$ está junto a usted sobre el asiento horizontal de su automóvil. El coeficiente de fricción estática entre el libro y el asiento es $0.650$, y el coeficiente de fricción cinética es $0.550$. Suponga que viaja a $72.0 \text{ km/h}$ y frena para detenerse con aceleración constante en una distancia de $30.0 \text{ m}$. Su libro de física permanece sobre el asiento y no resbala hacia adelante para caer al piso.
    \vspace*{6cm}

    \item[\textbf{27.}] Considere un camión grande que transporta una pesada carga, como vigas de acero. Un gran riesgo para el conductor es que la carga se deslice hacia adelante y golpee la cabina, si el camión se detiene repentinamente durante un accidente o al frenar. Suponga, por ejemplo, que una carga de $10\,000 \text{ kg}$ va colocada sobre un camión de $20\,000 \text{ kg}$ moviéndose a $12.0 \text{ m/s}$. Supongamos que la carga no está atada al camión, pero tiene un coeficiente de fricción de $0.500$ con la plataforma del camión. (a) Calcule la distancia mínima de parada para que la carga no se deslice hacia adelante en relación con el camión. (b) ¿Hay algún dato innecesario para la solución?
    \vspace*{6cm}

    \item[\textbf{28.}] Antes de 1960, se creía que el máximo coeficiente de fricción estática alcanzable para la llanta de un automóvil, sobre un camino, era $\mu_s = 1$. Alrededor de 1962, tres compañías, cada una por su lado, desarrollaron llantas de carreras con coeficientes de 1.6. Este problema muestra que las llantas han mejorado desde aquella ocasión. El intervalo de tiempo más corto para un automóvil con motor de pistones, inicialmente en reposo, para cubrir una distancia de un cuarto de milla es $4.43 \text{ s}$. (a) Suponga que las llantas traseras levantaron las delanteras del pavimento, como lo muestra la figura P5.28. ¿Qué valor mínimo de $\mu_s$ es necesario para lograr el tiempo récord? (b) Suponga que al conductor le fuera posible incrementar la potencia de su motor y mantener iguales otras cosas. ¿Cómo afectaría este cambio al intervalo de tiempo?
    \vspace*{6cm}

    \item[\textbf{29.}] Un objeto suspendido de $9.00 \text{ kg}$ se conecta, mediante una cuerda ligera inextensible sobre una polea ligera sin fricción, a un bloque de $5.00 \text{ kg}$ que se desliza sobre una mesa plana (figura P5.22). Si toma el coeficiente de fricción cinética como $0.200$, encuentre la tensión en la cuerda.
    \vspace*{6cm}
\end{itemize}

\newpage

\begin{itemize}
    \item[\textbf{30.}] En la figura P5.30, la persona pesa $170 \text{ lb}$. Visto de frente, cada muleta forma un ángulo de $22.0^\circ$ con la vertical. Las muletas soportan la mitad del peso del individuo. La otra mitad es soportada por las fuerzas verticales del suelo sobre los pies de la persona. Suponga que el hombre se mueve con velocidad constante y que la fuerza ejercida por el piso sobre las muletas actúa a lo largo de ellas, determine (a) el coeficiente de fricción más pequeño posible entre las muletas y el suelo, y (b) la magnitud de la fuerza de compresión en cada muleta.
    \vspace*{6cm}

    \item[\textbf{31.}] Tres objetos se conectan sobre una mesa, como se muestra en la figura P5.31. El coeficiente de fricción cinética entre el bloque de masa $m_2$ y la mesa es $0.350$. Los objetos tienen masas de $m_1 = 4.00 \text{ kg}$, $m_2 = 1.00 \text{ kg}$ y $m_3 = 2.00 \text{ kg}$, y las poleas no tienen fricción. (a) Dibuje un diagrama de cuerpo libre para cada objeto. (b) Determine la aceleración de cada objeto, incluyendo su dirección. (c) Determine las tensiones en las dos cuerdas. \textbf{¿Qué pasaría si?} (d) Si la mesa fuera lisa, ¿las tensiones aumentarían, disminuirían o quedarían intactas? Explique.
    \vspace*{8cm}

    \item[\textbf{32.}] Usted está trabajando como clasificador de cartas en el Servicio Postal de Estados Unidos. Las normas postales requieren que el calzado de sus empleados tenga un coeficiente de fricción estática de $0.5$ o más, sobre una superficie específica. Usted utiliza un calzado de atletismo cuyo coeficiente de fricción estática no conoce. Para determinar el coeficiente, usted imagina que hay una emergencia y empieza a correr por la habitación. Un compañero de trabajo le toma el tiempo, y encuentra que usted puede comenzar desde el reposo y moverse $4.23 \text{ m}$ en $1.20 \text{ s}$. Si intenta moverse más rápido que esto, sus pies resbalan. Suponiendo que su aceleración es constante, ¿su calzado califica para la regulación postal?
    \vspace*{6cm}

    \item[\textbf{33.}] Se le ha llamado como testigo experto para un juicio en el que un conductor ha sido acusado por exceso de velocidad, pero declara su inocencia. Afirma haber frenado para no chocar con la parte trasera de otro auto, pero la golpeó justo cuando se detuvo. Usted ha sido contratado por el fiscal para demostrar que el conductor de hecho iba a exceso de velocidad. Usted ha obtenido de la policía los datos siguientes: que las marcas de la patinada dejadas por el conductor son de $56.0 \text{ m}$ de largo y la carretera está a nivel. Los neumáticos coinciden con los del auto del conductor que se ha arrastrado por la misma carretera para determinar que el coeficiente de fricción cinética entre los neumáticos y la carretera es de $0.82$ en todos los puntos a lo largo de la marca de la patinada. El límite de velocidad en la carretera es de $35 \text{ mi/h}$. Redacte un argumento que pueda usar en la corte para demostrar que el conductor de hecho iba a exceso de velocidad.
    \vspace*{7cm}
\end{itemize}

\newpage

\begin{itemize}
    \item[\textbf{34.}] Un bloque de $3.00 \text{ kg}$ de masa es empujado contra una pared mediante una fuerza $\vec{P}$ que forma un ángulo $\theta = 50.0^\circ$ con la horizontal, como se muestra en la figura P5.34. El coeficiente de fricción estática entre el bloque y la pared es $0.250$. (a) Determine los valores posibles para la magnitud de $\vec{P}$ que permiten al bloque permanecer fijo. (b) Describa qué sucede si $|\vec{P}|$ tiene un valor mayor y qué ocurre si es más pequeño. (c) Repita los incisos (a) y (b) suponiendo que la fuerza forma un ángulo $\theta = 13.0^\circ$ con la horizontal.
    \vspace*{8cm}

    \item[\textbf{35.}] \textbf{Problema de repaso.} Un salmón real puede nadar a $3.58 \text{ m/s}$, y también puede saltar verticalmente, abandonando el agua con una rapidez de $6.26 \text{ m/s}$. Un salmón récord tiene $1.50 \text{ m}$ de longitud y una masa de $61.0 \text{ kg}$. Considere un pez que nada en línea recta hacia arriba bajo la superficie de un lago. La fuerza gravitacional ejercida sobre el pez es casi cancelada por la fuerza de flotación ejercida por el agua, como se estudiará en el capítulo 10. El pez experimenta una fuerza $P$ hacia arriba ejercida por el agua sobre las aletas de su cola y una fuerza hacia abajo debida a la fricción del fluido que se modela actuando sobre su extremo frontal. Suponga que la fuerza de fricción del fluido desaparece tan pronto como la cabeza del pez corta la superficie del agua, y asuma que la fuerza sobre su cola es constante. Modele la fuerza gravitacional que se elimina rápidamente cuando la mitad del pez está fuera del agua. Encuentre el valor de $P$.
    \vspace*{6cm}

    \item[\textbf{36.}] Un bloque de $5.00 \text{ kg}$ se coloca encima de otro de $10.0 \text{ kg}$ (figura P5.36). Una fuerza horizontal de $45.0 \text{ N}$ es aplicada al bloque de $10.0 \text{ kg}$, y el bloque de $5.0 \text{ kg}$ está amarrado a la pared. El coeficiente de fricción cinética entre todas las superficies en movimiento es $0.200$. (a) Dibuje un diagrama de cuerpo libre para cada bloque e identifique las fuerzas de acción-reacción entre los bloques. (b) Determine la tensión en la cuerda y la magnitud de la aceleración del bloque de $10.0 \text{ kg}$.
    \vspace*{7cm}
\end{itemize}

\newpage

\subsection*{PROBLEMAS ADICIONALES}

\begin{itemize}
    \item[\textbf{37.}] Un deslizador de aluminio negro flota sobre una película de aire en una pista de aire de aluminio a nivel. En esencia, el aluminio no siente fuerza en un campo magnético y la resistencia del aire es despreciable. Un imán intenso se une a lo alto del deslizador y forma una masa total de $240 \text{ g}$. Un trozo de chatarra de hierro unido a un tope en la pista atrae al imán con una fuerza de $0.823 \text{ N}$ cuando el hierro y el imán están separados $2.50 \text{ cm}$. (a) Encuentre la aceleración del deslizador en este instante. (b) La chatarra de hierro ahora se une a otro deslizador verde y forma una masa total de $120 \text{ g}$. Encuentre la aceleración de cada deslizador cuando se liberan simultáneamente a $2.50 \text{ cm}$ de separación.
    \vspace*{6cm}

    \item[\textbf{38.}] ¿Por qué no es posible la siguiente situación? Un libro se coloca sobre un plano inclinado en la superficie de la Tierra. El ángulo del plano con la horizontal es $60.0^\circ$. El coeficiente de fricción cinética entre el libro y el plano es $0.300$. Al tiempo $t = 0$, el libro se libera a partir del reposo. Entonces el libro se desliza una distancia de $1.00 \text{ m}$, medida a lo largo del plano, en un intervalo de tiempo de $0.483 \text{ s}$.
    \vspace*{6cm}

    \item[\textbf{39.}] Dos bloques de masas $m_1$ y $m_2$ se colocan sobre una tabla en contacto entre sí como se discute en el ejemplo 5.7 y se muestra en la figura 5.13a. El coeficiente de fricción cinética entre el bloque de la masa $m_1$ y la tabla es $\mu_1$, y entre el bloque de masa $m_2$ y la tabla es $\mu_2$. Se aplica una fuerza horizontal de magnitud $F$ al bloque de masa $m_1$. Deseamos encontrar $P$, la magnitud de la fuerza de contacto entre los bloques. (a) Dibuje los diagramas que demuestran las fuerzas para cada bloque. (b) ¿Cuál es la fuerza neta en el sistema de dos bloques? (c) ¿Cuál es la fuerza neta que actúa sobre $m_1$? (d) ¿Cuál es la fuerza neta que actúa sobre $m_2$? (e) Escriba la segunda ley de Newton en la dirección $x$ para cada bloque. (f) Resuelva las dos ecuaciones con dos incógnitas para la aceleración $a$ de los bloques en términos de las masas, la fuerza aplicada $F$, los coeficientes de fricción, y $g$. (g) Encuentre la magnitud $P$ de la fuerza de contacto entre los bloques en términos de las mismas cantidades.
    \vspace*{9cm}
\end{itemize}

\newpage

\begin{itemize}
    \item[\textbf{40.}] Un deslizador de $1.00 \text{ kg}$ sobre un riel de aire horizontal se jala con una cuerda a un ángulo $\theta$. La cuerda tensa pasa por una polea y está ligada a un cuerpo que cuelga y tiene una masa de $0.500 \text{ kg}$, como se muestra en la figura P5.40. (a) Muestre que la rapidez $v_x$ del deslizador y la rapidez $v_y$ del objeto que cuelga están relacionadas $v_x = u v_y$, donde $u = z(z^2 - h_0^2)^{-1/2}$. (b) El deslizador se libera a partir del reposo. Demuestre que en ese instante la aceleración $a_x$ del deslizador y la aceleración $a_y$ del objeto que cuelga están relacionadas por $a_x = u a_y$. (c) Encuentre la tensión en la cuerda en el instante que se libera el deslizador, para $h_0 = 80.0 \text{ cm}$ y $\theta = 30.0^\circ$.
    \vspace*{7cm}

    \item[\textbf{41.}] Un niño inventivo llamado Nick quiere alcanzar sin escalar una manzana que pende de un árbol. Sentado en una silla unida a una cuerda que pasa sobre una polea sin fricción (figura P5.41), Nick jala sobre el extremo suelto de la cuerda con tal fuerza que la balanza de resorte lee $250 \text{ N}$. El verdadero peso de Nick es $320 \text{ N}$ y la silla pesa $160 \text{ N}$. Los pies de Nick no tocan el suelo. (a) Dibuje un par de diagramas que muestren las fuerzas para Nick y la silla considerados como sistemas separados y otro diagrama para Nick y la silla considerados como un sistema. (b) Muestre que la aceleración del sistema es \textit{hacia arriba} y encuentre su magnitud. (c) Encuentre la fuerza que Nick ejerce sobre la silla.
    \vspace*{7cm}

    \item[\textbf{42.}] Una cuerda con masa $m_c$ está unida a un bloque de masa $m_b$, como en la figura P5.42. El bloque descansa sobre una superficie horizontal sin fricción. La cuerda no se estira. El extremo libre de la cuerda se jala hacia la derecha con una fuerza horizontal $\vec{F}$. (a) Dibuje diagramas de fuerzas para la cuerda y el bloque, observando que la tensión en la cuerda no es uniforme. (b) Encuentre la aceleración del sistema en términos de $m_c, m_b$ y $F$. (c) Encuentre la magnitud de la fuerza que la cuerda ejerce sobre el bloque. (d) ¿Qué pasa con la fuerza sobre el bloque conforme la masa de la cuerda se aproxima a cero? ¿Qué puede decirse acerca de la tensión en una cuerda \textit{ligera} uniendo a un par de objetos en movimiento?
    \vspace*{7cm}

    \item[\textbf{43.}] En el ejemplo 5.7 se empujaron dos bloques sobre una mesa. Suponga que tres bloques están en contacto mutuo sobre una superficie horizontal sin fricción, como se muestra en la figura P5.43. A $m_1$ se le aplica una fuerza horizontal $\vec{F}$. Tome $m_1 = 2.00 \text{ kg}$, $m_2 = 3.00 \text{ kg}$, $m_3 = 4.00 \text{ kg}$ y $F = 18.0 \text{ N}$. (a) Dibuje un diagrama de cuerpo libre para cada bloque. (b) Determine la aceleración de los bloques. (c) Encuentre la fuerza \textit{resultante} sobre cada bloque. (d) Encuentre las magnitudes de las fuerzas de contacto entre los bloques. (e) Usted trabaja en un proyecto de construcción. Un colaborador clava tablarroca en un lado de un separador ligero y usted está en el lado opuesto, proporcionando ``respaldo'' al apoyarse contra la pared con su espalda, empujando sobre ella. Cada golpe de martillo hace que su espalda sufra un pinchazo. El supervisor lo ayuda al poner un pesado bloque de madera entre la pared y su espalda. Use la situación analizada en los incisos (a), (b) y (c) como modelo, y explique cómo este cambio funciona para hacer su trabajo más confortable.
    \vspace*{9cm}
\end{itemize}

\newpage

\begin{itemize}
    \item[\textbf{44.}] En la situación descrita en el problema 41 y la figura P5.41, las masas de la cuerda, balanza y polea son despreciables. Los pies de Nick no tocan el suelo. (a) Suponga que Nick está momentáneamente en reposo cuando deja de jalar la cuerda hacia abajo y pasa el extremo de la cuerda a otro niño, de $440 \text{ N}$ de peso, que está de pie en el suelo junto a él. La cuerda no se rompe. Describa el movimiento resultante. (b) En vez de esto, suponga que Nick está momentáneamente en reposo cuando amarra el extremo de la cuerda a una saliente en forma de gancho resistente que se deriva del tronco del árbol. Explique por qué esta acción puede hacer que la cuerda se rompa.
    \vspace*{6cm}

    \item[\textbf{45.}] Una caja de peso $F_g$ es empujada mediante una fuerza $\vec{P}$ sobre un piso horizontal, como se muestra en la figura P5.45. El coeficiente de fricción estática es $\mu_s$, y $\vec{P}$ se dirige a un ángulo $\theta$ bajo la horizontal. (a) Demuestre que el valor mínimo de $P$ que moverá la caja está dado por
    $$ P = \frac{\mu_s F_g \sec\theta}{1 - \mu_s \tan\theta} $$
    (b) Encuentre la condición sobre $\theta$, en términos de $\mu_s$, tal que sea imposible el movimiento de la caja para cualquier valor de $P$.
    \vspace*{6cm}

    \item[\textbf{46.}] En la figura P5.46, la polea y la cuerda son ligeras, todas las superficies carecen de fricción, y la cuerda no se estira. (a) ¿Qué comparación puede hacerse entre las aceleraciones de los bloques 1 y 2? Explique su razonamiento. (b) La masa del bloque 2 es $1.30 \text{ kg}$. Encuentre su aceleración que depende de la masa $m_1$ del bloque 1. (c) \textbf{¿Qué pasaría si?} ¿Qué predice el resultado del inciso (b) si $m_1$ es mucho menor que $1.30 \text{ kg}$? (d) ¿Qué predice el resultado del inciso (b) si $m_1$ se aproxima a infinito? (e) En este último caso, ¿cuál es la tensión en la cuerda? (f) ¿Podría usted anticipar las respuestas de los incisos (c), (d) y (e) sin hacer primero el inciso (b)? Explique.
    \vspace*{8cm}
\end{itemize}

\newpage

\begin{itemize}
    \item[\textbf{47.}] Usted está trabajando como un testigo experto para la defensa de un capitán de un buque contenedor cuya nave se topó con un arrecife que rodea una isla caribeña. El capitán es acusado de dirigir la nave intencionalmente hacia el arrecife. Se ha descubierto la siguiente información que se ha presentado y los abogados de ambos bandos han estipulado que la información es correcta: el barco viajaba a $2.50 \text{ m/s}$ hacia el arrecife cuando una falla mecánica causó que el timón se atascara en la posición recta. En ese momento, la nave estaba a $900 \text{ m}$ del arrecife. El viento soplaba directamente hacia el arrecife y ejercía una fuerza constante de $9.00 \times 10^4 \text{ N}$ en el barco en dirección hacia el arrecife. La masa de la nave y su carga era de $5.50 \times 10^7 \text{ kg}$. Durante la preparación del juicio, el capitán afirma que sin el control de la dirección de viaje, la única opción que tenía era poner los motores en reversa a la potencia máxima, de tal manera que la fuerza total ejercida por la fuerza de arrastre de fricción del agua y la fuerza del agua en las hélices era de $1.25 \times 10^5 \text{ N}$ alejándose del arrecife. Con esta información, construya un argumento convincente de que el capitán no podía hacer nada en esta situación para evitar que el buque pegara contra el arrecife.
    \vspace*{7cm}

    \item[\textbf{48.}] Un cojín plano de masa $m$ se libera desde el reposo en lo alto de un edificio que tiene una altura $h$. Un viento que sopla a lo largo del lado del edificio ejerce una fuerza horizontal constante de magnitud $F$ sobre el cojín conforme cae, como se muestra en la figura P5.48. El aire no ejerce fuerza vertical. (a) Demuestre que la trayectoria del cojín es una línea recta. (b) ¿El cojín cae con velocidad constante? Explique. (c) Si $m = 1.20 \text{ kg}$, $h = 8.00 \text{ m}$ y $F = 2.40 \text{ N}$, ¿a qué distancia del edificio el cojín golpeará el nivel del suelo? \textbf{¿Qué pasaría si?} (d) Si el cojín se lanza hacia abajo con una rapidez distinta de cero, desde lo alto del edificio, ¿cuál será la forma de su trayectoria? Explique.
    \vspace*{8cm}

    \item[\textbf{49.}] Qué fuerza horizontal debe aplicarse a un gran bloque de masa $M$, que se muestra en la figura P5.49, para que los otros bloques permanezcan fijos respecto a $M$? Suponga que todas las superficies y la polea carecen de fricción. Note que la fuerza ejercida por la cuerda acelera a $m_2$.
    \vspace*{6cm}

    \item[\textbf{50.}] Un objeto de $8.40 \text{ kg}$ se desliza hacia abajo por un plano inclinado fijo sin fricción. Use una computadora para determinar y tabular (a) la fuerza normal que se ejerce sobre el objeto y (b) su aceleración para una serie de ángulos de inclinación (medidos desde la horizontal) que varían de $0^\circ$ a $90^\circ$ en incrementos de $5^\circ$. (c) Trace una gráfica de la fuerza normal y la aceleración como funciones del ángulo de inclinación. (d) En los casos límite de $0^\circ$ y $90^\circ$, ¿sus resultados son consistentes con el comportamiento conocido?
    \vspace*{7cm}
\end{itemize}

\newpage

\subsection*{PROBLEMAS DE DESAFÍO}

\begin{itemize}
    \item[\textbf{51.}] Un bloque de masa $2.20 \text{ kg}$ se acelera sobre una superficie rugosa mediante una cuerda ligera que pasa por una pequeña polea, como se indica en la figura P5.51. La tensión $T$ en la cuerda se mantiene a $10 \text{ N}$, y la polea está a $0.100 \text{ m}$ por arriba del bloque. El coeficiente de fricción cinética es $0.400$. (a) Determine la aceleración del bloque cuando $x = 0.400 \text{ m}$. (b) Describa el comportamiento general de la aceleración conforme el bloque se desliza desde una ubicación con $x$ grande hasta $x = 0$. (c) Encuentre el valor máximo de la aceleración y la posición $x$ en donde esto ocurre. (d) Encuentre el valor de $x$ para el cual la aceleración es cero.
    \vspace*{8cm}

    \item[\textbf{52.}] ¿Por qué no es posible la siguiente situación? Una tostadora de $1.30 \text{ kg}$ no está conectada. El coeficiente de fricción estática entre la tostadora y un mostrador horizontal es $0.350$. Para hacer que la tostadora comience a moverse, usted jala descuidadamente su cordón eléctrico. Desafortunadamente, el cordón se ha debilitado debido a acciones similares y se romperá si la tensión en él excede $4.00 \text{ N}$. Jalando el cordón en un ángulo particular, usted logra que la tostadora se mueva sin romper el cordón.
    \vspace*{6cm}

    \item[\textbf{53.}] Inicialmente, el sistema de objetos mostrado en la figura P5.49 se mantiene sin movimiento. La polea, todas las superficies y las ruedas carecen de fricción. Acepte que la fuerza $\vec{F}$ es cero y suponga que $m_1$ solo se puede mover verticalmente. En el instante después de liberar al sistema de objetos, encuentre (a) la tensión $T$ en la cuerda, (b) la aceleración de $m_2$, (c) la aceleración de $M$ y (d) la aceleración de $m_1$. (\textit{Nota}: La polea acelera junto con el sistema.)
    \vspace*{8cm}
\end{itemize}

\newpage

\begin{itemize}
    \item[\textbf{54.}] Un móvil se forma al soportar cuatro mariposas metálicas de igual masa $m$ de una cuerda de longitud $L$. Los puntos de soporte están igualmente espaciados una distancia $\ell$, como se ilustra en la figura P5.54. La cuerda forma un ángulo $\theta_1$ con el techo en cada punto final. La sección central de la cuerda es horizontal. (a) Encuentre la tensión en cada sección de cuerda en términos de $\theta_1, m$ y $g$. (b) Encuentre el ángulo $\theta_2$, en términos de $\theta_1$, que las secciones de cuerda entre las mariposas exteriores y las mariposas interiores forman con la horizontal. (c) Demuestre que la distancia $D$ entre los puntos extremos de la cuerda es
    $$ D = \frac{L}{5} \left\{ 2 \cos\theta_1 + 2\cos[\tan^{-1}(\tfrac{1}{2}\tan\theta_1)] + 1 \right\} $$
    \vspace*{8cm}

    \item[\textbf{55.}] En la figura P5.55, el plano inclinado tiene masa $M$ y está clavado a la mesa horizontal. El bloque de masa $m$ se coloca cerca del fondo del plano y es liberado con un empujón para deslizarlo hacia arriba. El bloque se coloca cerca de lo alto del plano, como se muestra en la figura, y después se desliza hacia abajo, siempre sin fricción. Encuentre la fuerza que la mesa ejerce sobre el plano durante este movimiento en términos de $m, M, g$ y $\theta$.
    \vspace*{7cm}
\end{itemize}

\end{document}
